\subsection{A common structure across the four mitigations}
\label{sec:mitigations}

The mitigations adopted in this paper are not independent engineering fixes.
Each was arrived at separately, from direct numerical evidence, a stalled
residual, a spiked row norm, a badly conditioned submap, and the common
structure was recognised only afterwards. Recording it is worth a paragraph,
because it suggests where to look first when the architecture is carried into
a solver where the specific failures differ.

Three of the four resolve a tension by separating the two things in conflict
rather than by compromising between them. The transition closure of
Section~\ref{sec:transition} separates them in \emph{time}: the conflict is
between a smooth, Newton-tractable residual and an accurate transition model,
and it is resolved by phase, forcing transition during the inverse iterations
and releasing it afterwards, with a verification solve confirming that nothing
was lost. Nothing is traded away in the end state; the kink is removed exactly
when Newton needs it removed. The prescribed leading edge separates them in
\emph{space}: the conflict is between design freedom, invert everywhere, and
row conditioning, since the residual rows near the nose blow up, and it is
resolved by prescribing the first few percent of chord geometrically and
inverting only aft of it. That does cost design freedom, and
Section~\ref{sec:limitations} records the cost rather than explaining it away.
The station-restriction rule of Section~\ref{sec:le-identifiability} separates
them along the same axis, one level down: having removed a region from the
design space, the sampling must be removed from it too.

The remaining two resolve their tension by improving the quality of what is
already there rather than by adding machinery. For the basis-conditioning
mode, the obvious fix, Tikhonov regularisation, turns a determined solve back
into a least-squares problem, which is the optimisation this architecture
exists to avoid; the discipline adopted instead is to keep the basis order
inside its well-conditioned regime, or to orthogonalise the basis, rather than
regularise a badly-conditioned one. For station selection, adding more target
equations would over-determine the square system, and accepting whatever
stations are convenient is what produced the wrong root of
Section~\ref{sec:ablations}; the resolution reuses an operator the pre-solve
had already built to choose which $M$ rows are maximally informative. Both are
instances of the same instinct: when a determined formulation is the point,
the fix must not add degrees of freedom.

\subsection{Threats to validity}
\label{sec:validity}

The most important limitation of this study is structural and is stated first
because it conditions how every accuracy number in
Section~\ref{sec:results} should be read. No result reported here recovers a
geometry from a target that was produced independently of the method. In every
case, the self-consistency test, the ablation matrix and both panels, the
target pressure distribution is generated by the same solver, at the same
paneling, from a coefficient vector that is exactly representable in the same
CST basis at the same order. Recovery errors at the $10^{-11}$ level are
therefore measurements of \emph{identifiability and conditioning}, showing
that the sampled square system has the generating design as its unique
reachable root and that the Newton iteration finds it, and not measurements of
accuracy against independent data. This is deliberate: a self-generated target
is the only construction under which a failure isolates cleanly to one of the
Table~\ref{tab:fm} failure modes rather than to target infeasibility, which is
what makes the test falsifiable in the first place. But it means the reported
precision does not transfer to a target drawn from experiment, from a finer
discretisation, or from a different solver. The nearest thing to evidence outside this regime is the
user-supplied-target path described in Section~\ref{sec:artifacts}, which
recovers shape to a fraction of a millichord while leaving coefficients to
differ at the $10^{-2}$ level. Even that solve was run against a
self-generated target, and its coefficient gap is attributable to the
inviscid pre-solve meeting a viscous target and to the basis conditioning of
Section~\ref{sec:nsweep} rather than to target independence, so it bounds
nothing. A genuinely independent target has not been attempted, and doing so
is the first item of future work.

Two further gaps in the evidence follow from the same source. The
realisability guard of Section~\ref{sec:presolve} is never exercised on an
unrealisable target: measured realisability on the reference configuration is
$3.90\x10^{-4}$ against a threshold of $0.05$, because a self-generated target
is realisable by construction, so the guard's discriminating power is asserted
rather than demonstrated. And there is no robustness study: real inverse
targets are noisy, over- or under-specified, and nothing here tests target
perturbation or a target that is only nearly realisable. Given that the
identifiability mechanism of Section~\ref{sec:ablations} is a statement about
how much information the sampled rows carry, a sweep of recovery error against
target-noise amplitude is the single cheapest experiment that would sharpen
this paper's central finding, and it has not been run.

The discretisation is also fixed throughout. Every result uses the same panel
count, and for a method whose target rows attach to surface stations and whose
flow-block Jacobian is a finite difference over geometry, sensitivity to panel
count is a first-order question that this study does not answer.

\subsection{Scope limits}
\label{sec:limitations}

The ablation evidence is single-airfoil and single-run. The order sweep, the
station-selection comparison, the remaining single-factor ablations of
Section~\ref{sec:other-ablations} and the cost comparison of
Section~\ref{sec:cost} are all on NACA 2412 at a fixed seed, one run per
configuration, with no confidence interval reported for any of them. That caveat does not extend to the recovery claim itself, which was
upgraded to panel-level evidence across both the 20-section NACA panel and the
117-section UIUC panel (Sections~\ref{sec:h1-panel}
and~\ref{sec:h1-uiuc-panel}); but the ablation-level findings remain
single-airfoil and are not extended to either panel here.

The panel evidence carries a selection effect, quantified in
Section~\ref{sec:h1-uiuc-panel}. Twenty-nine percent of the UIUC panel is
excluded, overwhelmingly because the host analysis solver cannot converge a
target for the section at all, and the excluded sections are systematically
thinner than the converged ones. The exclusion arises before any inverse
problem is posed and is therefore a property of the host solver rather than of
the formulation, but the accuracy claim is conditioned on it. Three further
sections fail the host solver's wake-direction assertion for a cause unrelated
to the sharp-trailing-edge case that an earlier correction addressed; their
trailing-edge gap is already above the minimum, and this remains an open item.

The flow regime is narrow. All results are incompressible: the configuration
sets free-stream Mach number to zero, so the solver's Kármán--Tsien
compressibility correction, which would in any case be valid only below the
critical Mach number, is never exercised. Compressible and transonic behaviour
is untested, as is anything cascade-specific, since there is no periodicity,
no pitch or stagger, and no streamtube contraction. All runs share a single
operating condition, one Reynolds number and one incidence, so no claim here
generalises across flow regimes. The hypothesis under test is a claim about
solver \emph{architecture} rather than about turbine aerodynamics, and is
answered as well on an isolated airfoil as on a cascade; but the regimes the
architecture was ultimately built for are genuinely untested.

Two limits internal to the formulation were flagged in
Section~\ref{sec:contribution} and are repeated here because they bound the
contribution rather than the evidence. The constraint rows, which are half of
the novelty claim, are implemented and unit-tested but not exercised by any
measured result, so wherever an exact analytic constraint-row Jacobian is
offered as an advantage over the nested baseline, it is an advantage of the
formulation and not one this study has measured, in a comparison where neither
side carried an active constraint. And the flow-block Jacobian is obtained by
finite differences rather than from the exact constant geometric sensitivity,
because the host solver exposes no complete geometric sensitivity of its own
residual; the exact sensitivity is used where it can be, in the constraint
rows and the pre-solve, and Section~\ref{sec:h3} measures what the
finite-difference substitute costs in convergence rate.

Finally, the novelty claim itself rests on a negative result established by
open literature search rather than by an exhaustive search of the aerospace
society digital collections. No published work is known that couples a
Newton-coupled inverse solver directly to CST coefficients as Newton unknowns,
as opposed to CST as a post-hoc smoother, but that is a statement about what
was found rather than about what exists.

\subsection{Why the pre-registered cost threshold was missed}
\label{sec:outdated-baseline}

The pre-registered hundredfold claim is not simply missed by the measured
$3.1$ to $8.1\x$ in counted solves, or by the $3.4\x$ in wall-clock time. The
two fair-paired controls of Section~\ref{sec:cost} give a plausible account of
\emph{why} it lands where it does. A Levenberg--Marquardt solve with all three
tolerances at $10^{-8}$ and per-variable scaling taken from the initial guess,
over the same 16 well-conditioned free coefficients, is a genuinely strong
comparator. It converges the cold-start case in six outer iterations, which is
nothing like the surrogate-in-the-loop nested-optimisation baselines that a
hundredfold figure typically describes in the older aerodynamic
shape-optimisation literature the original framing drew on. The pre-registered
threshold plausibly encodes an outdated class of nested baseline. Measuring
against a strong modern gradient-based one and reporting the resulting ratio
directly, in three currencies rather than the most favourable one, is a more
informative comparison than either confirming or missing the original number.
